\scanpage{54}
d'un objet (espace) simplicial --- \uline{au dessus de}
$\Gamma(\varnothing)=\Sigma^{<}(0)$. Ici il s'agit d'un espace
simplicial relatif, et les conditions $b_1,b_2,b_3$
concernent cette structure --- elles disent
essentiellement que cette structure est associée
à une stratification cylindrique locale de
$\Sigma^{<}(0)=\Sigma$, et ``ordonnée'' (i.e. l'ens des
strates de codim $1$ passant par un $x\in\Sigma$ sont
\deleted{\unsure{est}} ordonnés, et cette relation d'ordre varie
continûment avec $x$ en un sens évident).
\margin{NB $\Gamma(n)=\unsure{\Sigma^{*}(n+1)|\Sigma^{<}(n+1)}$}
Les \deleted{$\Sigma^{<}(d)$ s'ident} $\Gamma(d)$\inserted{$=\Sigma^{<}(d+1)$} s'identifient alors
aux \deleted{parties} espaces des \inserted{(germes des)} strates de codim
$d+1$, \deleted{les} \inserted{\unsure{qu'on}} identifiés \deleted{grâce à la structure
d'ordre)} aux ens de germes de strates
de codim $1$ \deleted{\unsure{de}} \inserted{\unsure{en un même}} $x$. Les applic\unsure{ns}
semi-simpliciales entre les $\Gamma(d)$ proviennent
alors de la structure d'ordre --- \deleted{\unsure{on y}}
peut reconstituer les espaces de drapeaux
associés à la stratification, à l'aide des
\deleted{\unsure{ces}} $\Gamma(d)$, grâce à cette structure ordonnée.
En résumé, la donnée de $\Gamma:d\unsure{\mid\! -\!|\longmapsto}\Gamma(d)$,
\deleted{\unsure{i.e. de}} (i.e. de $\Sigma^{<}$ avec les appl\unsure{ns} \deleted{\unsure{semi}} simpliciales
cofinales) équivaut à celle d'un $\Sigma=\Sigma^{<}(0)$,
\mspage{53}
